\section{Background and related work}\label{sec:background}

\subsection{Decoy-state BB84 and finite-key security}
Weak coherent pulses carry a Poisson photon number, so multi-photon pulses leak information through photon-number splitting \cite{brassard2000}. The decoy-state method \cite{hwang2003,lo2005,wang2005} sends several intensities at random; since an eavesdropper sees only the photon number, the yields $y_n$ and error rates $e_n$ cannot depend on the intensity, and linear programming over the observed gains bounds the single-photon yield and its phase error. Practical versions use one signal, one weak decoy and one vacuum-like intensity \cite{ma2005}. Lim \etal \cite{lim2014} give closed-form finite-key bounds for this three-intensity protocol, using Hoeffding-type deviations for the decoy estimates and a sampling bound for the phase error. Hoeffding's inequality presumes independent trials, so these decoy bounds are rigorous for attacks that act independently on each pulse; coherent attacks call for martingale-type concentration bounds, for example Kato's inequality as used in \cite{zhang2017}. Tighter analyses exist \cite{tomamichel2012}, but the closed form suits a benchmark that evaluates thousands of blocks, and the attacks studied here are pulse-wise independent. The bounds also assume exactly known intensities. Intensity fluctuations have been treated rigorously, for random errors of known range \cite{wang2007}, for imperfect sources in the finite-key regime \cite{mizutani2015}, and with improved bounds for practical decoy-state systems \cite{zhang2017}. We do not use these analyses; we substitute the block mean of a per-window intensity monitor into the bounds of \cite{lim2014}, which is a heuristic, and check its effect against the true single-photon counts that the simulation records (Section~\ref{sec:results}).

\subsection{Quantum hacking}
Implementation attacks exploit the gap between device and model \cite{scarani2009,xu2020,pirandola2020}. Time-shift attacks use detector efficiency mismatch \cite{zhao2008}; Trojan-horse attacks read the modulator settings through back-reflected light \cite{gisin2006,lucamarini2015}; blinding attacks remove the quantum behaviour of the detectors altogether \cite{makarov2005,lydersen2010,gerhardt2011}, including superconducting nanowire detectors \cite{lydersen2011}; attacks on the source change the mean photon number that the decoy analysis assumes \cite{sajeed2015,huang2019}, and an attacker who can influence device calibration can open a detector side channel \cite{jain2011}. Countermeasures act on the hardware or on richer detection statistics: photocurrent monitoring \cite{yuan2010}, random detector efficiency \cite{lim2015}, whose assumptions were later broken in a commercial system \cite{huang2016}, variable attenuation combined with detection and error statistics \cite{qian2019}, real-time monitoring of inter-detection times \cite{dasilva2012}, background-noise cancellation for low-noise detectors \cite{lee2016}, a test on two-detector coincidences that bounds Eve's information \cite{gras2021}, best-practice criteria for self-differencing detectors \cite{koehlersidki2018}, and measurement-device-independent QKD \cite{lo2012}, which removes detector side channels by design. Several of these are statistical tests, but on observables that the window counts studied here do not contain.

\subsection{Telemetry-based attack detection}
Machine-learning detectors were first proposed for continuous-variable QKD, where the shot-noise and excess-noise estimates form natural features \cite{mao2020,liao2022}. For discrete-variable QKD, Al-Kuwari \etal \cite{alkuwari2026} train hybrid quantum-classical LSTMs on 10\,329 simulated decoy-state samples over eight balanced scenarios. Isha \etal \cite{isha2026} classify a balanced set of 24\,000 simulated BB84 sessions with gradient-boosted trees and report a false-negative rate of 0.0198 against 0.8477 for an 11\% QBER threshold; they also compare with a CUSUM detector, which reaches 78.99\% accuracy but at a high false-positive rate at the chosen threshold. Mohamed and Al-Kuwari \cite{mohamed2026} go furthest: they pose detection as a game against an adaptive adversary with bounded hardware, use block-level telemetry that includes decoy-state residuals, timing moments and detector imbalance, apply domain randomisation, accumulate evidence across blocks with a CUSUM, and weight missed detections by the loss of finite-key secret fraction computed with three-intensity decoy estimators and entropy accumulation. Machine learning has also been used on the defender's side for parameter optimisation \cite{wang2019}.

Against this background the present study has a narrower aim. It does not propose a better classifier. It differs from \cite{mohamed2026}, the closest work, in four respects: decisions are made per 10-second window rather than per post-processing block; benign data are 24-hour streams in which loss, efficiency, background, misalignment and intensities drift slowly and in a correlated way, rather than training blocks whose parameters are randomised independently; every attacked window is compared with the same window of an honest copy of the channel; and each block's key is audited against a ground-truth count of the bits Eve knows rather than against the bound. With these tools it asks what count telemetry can reveal at all when the attacker reproduces the monitored statistics, how often a detector alarms on a healthy drifting link, and whether the finite-key analysis would have released a key the attacker knows.

\subsection{Drift and false alarms}
Field links drift. Polarisation and fibre length change over tens of minutes, and long-running systems close feedback loops on the QBER and the count rates to compensate \cite{dixon2010}; network trials have since run QKD layers in field and production settings, the SwissQuantum network for more than a year \cite{stucki2011} and the Tokyo network as a field testbed \cite{sasaki2011}, so a monitor must cope with environmental variation over days rather than minutes. In the stream-learning literature this is concept drift \cite{gama2014}: a model fitted to one regime degrades when the input distribution moves. For a monitoring system the visible symptom is the false-alarm rate. Sequential change detection, from Page's CUSUM onwards \cite{page1954}, treats false alarms as a design constraint rather than a by-product; that view drives the design of the statistic in Section~\ref{sec:dcs}.
